\subsection{Divisibility of rational integers}

As was pointed out in section 1, the ring Z of rational integers is a principal ideal domain; its field of fractions is Q. The multiplicative group U of invertible elements of Z has two elements 1 and -1. The group \(Q_{+}^{*}\) of rational numbers \textgreater0 contains precisely one element from each class of associate elements of Q; it is thus isomorphic to the multiplicative group \(P^{*} = Q^{*}/U\) of principal fractional ideals of Q, with which it is usually identified. In particular, whenever gcd or lcm is used in the field Q (with respect to the ring Z), it is understood that these are elements \(\geqslant 0\) ; this convention allows us to speak of the gcd and the lcm of a family of rational numbers.

The irreducible integers \textgreater0 in Z are precisely those we have called prime numbers (I, p.~50) (they are sometimes called rational prime numbers); every irreducible element of Z is thus of the form p or -p, where p is a prime number, and the set P of prime numbers is a system of representatives of irreducible elements of Z.

PROPOSITION 5. --- The set of prime numbers is infinite.

Indeed, given an arbitrary finite family \((p_{i})\) \((1 \leqslant i \leqslant n)\) of distinct prime numbers, any prime divisor q of the number \(\left(\prod_{i=1}^{n} p_{i}\right) + 1\) (which is \textgreater1) is distinct from all the \(p_{i}\) , for otherwise it would divide 1.

